% GENERATED FILE. Edit canonical lessons, then run npm run generate:books.
% canonical-source-hash: fnv1a64:ba6bf508f205b9c5
% canonical-lessons: AR-C491-fakhur, AR-C491-iftakhara, AR-C491-imtinan, AR-C491-mumtann, AR-C491-ajaba

\chapter{Pride and thanks}
\label{ch:ar-pride-and-thanks}
\hlchaptermodality{491}

\begin{chapteropening}
\textbf{By the end of this chapter:} \emph{I can say that someone is proud or grateful, talk about gratitude, and say that something impressed me.}

Complete the last lesson of chapter 491: I can say that someone is proud or grateful, talk about gratitude, and say that something impressed me.
\end{chapteropening}

\section[\texorpdfstring{\ar{فخور} (fakhūr)}{فخور (fakhūr)}]{\ar{فخور} (fakhūr) — proud}

\label{lesson:AR-C491-fakhur}



\begin{quote}
Before the new one: say the Arabic for to be surprised, then the Arabic for admiration.
\end{quote}

\subsection*{You'll want to know: \ar{فخور}}
\textbf{\ar{فخور}} — \emph{fakhūr} — \textquotedblleft{}proud\textquotedblright{}.

\begin{quote}
\ar{الأم} \ar{فخورة} \ar{بابنها} \ar{الطبيب}.
\end{quote}

\begin{quote}
\emph{al-umm fakhūra bi-ibnihā aṭ-ṭabīb.}
\end{quote}

\begin{quote}
The mother is proud of her son the doctor.
\end{quote}

It takes \ar{بـ}: \ar{فخور} \ar{بـ}, proud of.

\begin{cousinweb}
\textbf{\ar{فخور}} (\emph{fakhūr}) is built on the native Arabic root \emph{f-kh-r}, \textquotedblleft{}to boast, to glory in\textquotedblright{} (Lane). Fakhūr is a pattern for a lasting trait: one who takes pride.
\end{cousinweb}

\begin{grammarlens}[title={What you've built}]
One more word for describing what you lived and felt.
\end{grammarlens}

\subsection*{Your turn}
\begin{itemize}
\raggedright
  \item \emph{Say it:} \emph{fakhūr}
  \item \emph{Say it:} \emph{fakhūr}, once more
  \item \emph{Say it:} \emph{al-umm fakhūra bi-ibnihā aṭ-ṭabīb.} — The mother is proud of her son the doctor.
  \item \emph{Recall:} say the Arabic for to be surprised, then the Arabic for admiration, then say \emph{fakhūr} again
\end{itemize}

\begin{tcolorbox}[breakable,colback=teal!4,colframe=teal!35!black,title={Before you move on}]
What does \ar{فخور} mean? (\textquotedblleft{}Proud\textquotedblright{}.) Say it once more.
\end{tcolorbox}
\section[\texorpdfstring{\ar{افتخر} (iftakhara)}{افتخر (iftakhara)}]{\ar{افتخر} (iftakhara) — to take pride (in)}

\label{lesson:AR-C491-iftakhara}



\begin{quote}
Before the new one: say the Arabic for admiration, then the Arabic for proud.
\end{quote}

\subsection*{You'll want to know: \ar{افتخر}}
\textbf{\ar{افتخر}} — \emph{iftakhara} — \textquotedblleft{}to take pride (in)\textquotedblright{}.

\begin{quote}
\ar{يفتخر} \ar{جدي} \ar{بحديقته} \ar{الجميلة}.
\end{quote}

\begin{quote}
\emph{yaftakhiru jaddī bi-ḥadīqatihi al-jamīla.}
\end{quote}

\begin{quote}
My grandfather takes pride in his beautiful garden.
\end{quote}

Present: \ar{يفتخر}, \emph{yaftakhir}.

It takes \ar{بـ}, like \ar{فخور}.

\begin{cousinweb}
\textbf{\ar{افتخر}} (\emph{iftakhara}) is built on the native Arabic root \emph{f-kh-r}, \textquotedblleft{}to boast, to glory in\textquotedblright{} (Lane). It is Form VIII of the root.
\end{cousinweb}

\begin{grammarlens}[title={What you've built}]
One more word for describing what you lived and felt.
\end{grammarlens}

\subsection*{Your turn}
\begin{itemize}
\raggedright
  \item \emph{Say it:} \emph{iftakhara}
  \item \emph{Say it:} \emph{iftakhara}, once more
  \item \emph{Say it:} \emph{yaftakhiru jaddī bi-ḥadīqatihi al-jamīla.} — My grandfather takes pride in his beautiful garden.
  \item \emph{Recall:} say the Arabic for admiration, then the Arabic for proud, then say \emph{iftakhara} again
\end{itemize}

\begin{tcolorbox}[breakable,colback=teal!4,colframe=teal!35!black,title={Before you move on}]
What does \ar{افتخر} mean? (\textquotedblleft{}To take pride (in)\textquotedblright{}.) Say it once more.
\end{tcolorbox}
\section[\texorpdfstring{\ar{امتنان} (imtinān)}{امتنان (imtinān)}]{\ar{امتنان} (imtinān) — gratitude}

\label{lesson:AR-C491-imtinan}



\begin{quote}
Before the new one: say the Arabic for proud, then the Arabic for to take pride.
\end{quote}

\subsection*{You'll want to know: \ar{امتنان}}
\textbf{\ar{امتنان}} — \emph{imtinān} — \textquotedblleft{}gratitude\textquotedblright{}.

\begin{quote}
\ar{كتبت} \ar{رسالة} \ar{امتنان} \ar{إلى} \ar{الممرضة}.
\end{quote}

\begin{quote}
\emph{katabtu risālat imtinān ilā al-mumarriḍa.}
\end{quote}

\begin{quote}
I wrote a letter of thanks to the nurse.
\end{quote}

\begin{cousinweb}
\textbf{\ar{امتنان}} (\emph{imtinān}) is built on the native Arabic root \emph{m-n-n}, \textquotedblleft{}to grant a favour\textquotedblright{} (Lane). Imtinān is the feeling of one who has been granted a favour.
\end{cousinweb}

\begin{grammarlens}[title={What you've built}]
One more word for describing what you lived and felt.
\end{grammarlens}

\subsection*{Your turn}
\begin{itemize}
\raggedright
  \item \emph{Say it:} \emph{imtinān}
  \item \emph{Say it:} \emph{imtinān}, once more
  \item \emph{Say it:} \emph{katabtu risālat imtinān ilā al-mumarriḍa.} — I wrote a letter of thanks to the nurse.
  \item \emph{Recall:} say the Arabic for proud, then the Arabic for to take pride, then say \emph{imtinān} again
\end{itemize}

\begin{tcolorbox}[breakable,colback=teal!4,colframe=teal!35!black,title={Before you move on}]
What does \ar{امتنان} mean? (\textquotedblleft{}Gratitude\textquotedblright{}.) Say it once more.
\end{tcolorbox}
\section[\texorpdfstring{\ar{ممتن} (mumtann)}{ممتن (mumtann)}]{\ar{ممتن} (mumtann) — grateful}

\label{lesson:AR-C491-mumtann}



\begin{quote}
Before the new one: say the Arabic for to take pride, then the Arabic for gratitude.
\end{quote}

\subsection*{You'll want to know: \ar{ممتن}}
\textbf{\ar{ممتن}} — \emph{mumtann} — \textquotedblleft{}grateful\textquotedblright{}.

\begin{quote}
\ar{أنا} \ar{ممتن} \ar{لك} \ar{على} \ar{مساعدتك}.
\end{quote}

\begin{quote}
\emph{anā mumtann laka ʿalā musāʿadatik.}
\end{quote}

\begin{quote}
I am grateful to you for your help.
\end{quote}

\ar{ممتن} \ar{لـ} someone \ar{على} something.

\begin{cousinweb}
\textbf{\ar{ممتن}} (\emph{mumtann}) is built on the native Arabic root \emph{m-n-n}, \textquotedblleft{}to grant a favour\textquotedblright{} (Lane). Mumtann is the participle that goes with imtinān, gratitude.
\end{cousinweb}

\begin{grammarlens}[title={What you've built}]
One more word for describing what you lived and felt.
\end{grammarlens}

\subsection*{Your turn}
\begin{itemize}
\raggedright
  \item \emph{Say it:} \emph{mumtann}
  \item \emph{Say it:} \emph{mumtann}, once more
  \item \emph{Say it:} \emph{anā mumtann laka ʿalā musāʿadatik.} — I am grateful to you for your help.
  \item \emph{Recall:} say the Arabic for to take pride, then the Arabic for gratitude, then say \emph{mumtann} again
\end{itemize}

\begin{tcolorbox}[breakable,colback=teal!4,colframe=teal!35!black,title={Before you move on}]
What does \ar{ممتن} mean? (\textquotedblleft{}Grateful\textquotedblright{}.) Say it once more.
\end{tcolorbox}
\section[\texorpdfstring{\ar{أعجب} (aʿjaba)}{أعجب (aʿjaba)}]{\ar{أعجب} (aʿjaba) — to impress, to please}

\label{lesson:AR-C491-ajaba}



\begin{quote}
Before the new one: say the Arabic for gratitude, then the Arabic for grateful.
\end{quote}

\subsection*{You'll want to know: \ar{أعجب}}
\textbf{\ar{أعجب}} — \emph{aʿjaba} — \textquotedblleft{}to impress, to please\textquotedblright{}.

\begin{quote}
\ar{أعجبني} \ar{الفيلم} \ar{كثيرا}.
\end{quote}

\begin{quote}
\emph{aʿjabanī al-fīlm kathīran.}
\end{quote}

\begin{quote}
I liked the film a lot.
\end{quote}

Present: \ar{يعجب}, \emph{yuʿjib}.

The one pleased is the object: \ar{أعجبني} \ar{الفيلم}, the film pleased me.

\begin{cousinweb}
\textbf{\ar{أعجب}} (\emph{aʿjaba}) is built on the native Arabic root \emph{ʿ-j-b}, \textquotedblleft{}to wonder, to marvel\textquotedblright{} (Lane). It is Form IV: to make someone marvel. Iʿjāb, admiration, is its verbal noun.
\end{cousinweb}

\begin{grammarlens}[title={What you've built}]
One more word for describing what you lived and felt.
\end{grammarlens}

\subsection*{Your turn}
\begin{itemize}
\raggedright
  \item \emph{Say it:} \emph{aʿjaba}
  \item \emph{Say it:} \emph{aʿjaba}, once more
  \item \emph{Say it:} \emph{aʿjabanī al-fīlm kathīran.} — I liked the film a lot.
  \item \emph{Recall:} say the Arabic for gratitude, then the Arabic for grateful, then say \emph{aʿjaba} again
\end{itemize}

\begin{tcolorbox}[breakable,colback=teal!4,colframe=teal!35!black,title={Before you move on}]
What does \ar{أعجب} mean? (\textquotedblleft{}To impress, to please\textquotedblright{}.) Say it once more.
\end{tcolorbox}
